\section*{अध्याय \ref{ch:b2:field-lab-methods-2} --- क्षेत्र और प्रयोगशाला विधियाँ II}
% numbers from the Book 4 second-writer scratch calculator calc/ch31.py

\begin{solution}{exo:b2:field-lab-methods-2:1}
(क) (i) \omterm{def:b2:field-lab-methods-2:textures}{ओफ़िटिक बनावट}; (ii) \omterm{def:b2:field-lab-methods-2:textures}{ग्रैनोब्लास्टिक बनावट}; (iii) एक \omterm{def:b2:field-lab-methods-2:porphyroblast}{पॉइकिलोब्लास्ट}, अंतर्वेशों से भरा एक
\omterm{def:b2:field-lab-methods-2:porphyroblast}{पॉर्फ़िरोब्लास्ट}; (iv) एक \omterm{def:b2:field-lab-methods-2:porphyroblast}{कोरोना}, टूट रहे किसी गार्नेट के चारों ओर एक अभिक्रिया-घेरा। (ख)
\eqref{eq:b2:field-lab-methods-2:lami} से, $\gamma_1/\gamma_2 = \sin\ang{104}/\sin\ang{128} =
0.9703/0.7880 = 1.23$: वह सीमा अन्य दोनों से \qty{23}{\percent} अधिक ज़ोर से खींचती है।
\end{solution}

\begin{solution}{exo:b2:field-lab-methods-2:2}
(क) $Q + A + P = 80$: Q $= 100 \times 25/80 = 31.3$, A $= 18.8$, P $= 50.0$। (ख) P/(A+P) $=
40/55 = \qty{72.7}{\percent}$, 65 और 90 के बीच, 20 और 60 के बीच Q के साथ: एक ग्रैनोडायोराइट।
(ग) $M' = 12 + 8 = \qty{20}{\percent}$, 35 से नीचे: ल्यूकोक्रेटी।
\end{solution}

\begin{solution}{exo:b2:field-lab-methods-2:3}
(क) $\sigma = \sqrt{0.12 \times 0.88/800} = 0.0115$, अर्थात् $12.0 \pm 1.1$ प्रतिशत-अंक,
$0.0115/0.12 = \qty{9.6}{\percent}$ की सापेक्ष त्रुटि। (ख) चार गुना बिंदु उसे आधा कर देते हैं: $0.0057$।
\end{solution}

\begin{solution}{exo:b2:field-lab-methods-2:4}
(क) हानि \qty{0.0288}{g} है, और $0.0288/0.8000 = \qty{3.60}{\percent}$।
(ख) गुटिका में नमूने का द्रव्यमान-अंश $0.7712/8.7712 = 0.0879$ है।
(ग) प्रज्वलित चूर्ण: $5.20/0.08792 = \qty{59.1}{\percent}$; मूल चूर्ण:
$59.14 \times (1 - 0.0360) = \qty{57.0}{\percent}$।
\end{solution}

\begin{solution}{exo:b2:field-lab-methods-2:5}
(क) $\sigma = 0.10 \times 0.02 = 0.002$ और $n = 0.02 \times 0.98/0.002^2 = 4900$ बिंदु।
(ख) $\sigma = 0.03$ और $n = 0.30 \times 0.70/0.03^2 = 233.3$, अतः 234 बिंदु। (ग) किसी निश्चित सापेक्ष
त्रुटि के लिए, $n = (1-p)/(p\,\varepsilon^2)$ $1/p$ की तरह बढ़ता है: सहायक खनिज 21 गुना अधिक
बिंदु माँगता है, जितने कोई काट प्रायः देती है उससे अधिक।
\end{solution}

\begin{solution}{exo:b2:field-lab-methods-2:6}
(क) RGM-1: $72.89/73.45 = \qty{99.2}{\percent}$; BCR-2: $54.17/54.1 = \qty{100.1}{\percent}$।
(ख) RGM-1: $0.19/\sqrt{20} = 0.042$, और $-0.56/0.042 = -13.2$; BCR-2: $0.13/\sqrt{20} = 0.029$, और
$0.07/0.029 = 2.4$। (ग) RGM-1 में पूर्वाग्रह है: उसका माध्य अनुशंसित मान से अपने ही मानक विचलन के 13 गुना से
चूकता है, संयोग की अनुमति से कहीं अधिक (यदि अनुशंसित मान सही है)। BCR-2 उस सीमा पर है जो
संयोग देता है। RGM-1 कम परिशुद्ध नहीं है: उसका सापेक्ष मानक विचलन,
$0.19/72.89 = \qty{0.26}{\percent}$, BCR-2 के \qty{0.24}{\percent} के निकट है। समस्या
यथार्थता की है, परिशुद्धता की नहीं।
\end{solution}

\begin{solution}{exo:b2:field-lab-methods-2:7}
(क) $M = 39.098 + 3 \times 24.305 + 26.982 + 3 \times 28.085 + 12 \times 15.999 + 2 \times 1.008 =
\qty{417.25}{g.mol^{-1}}$; दोनों OH मिलकर एक \ce{H2O} ($18.015$) बनाते हैं, \qty{4.32}{\percent}।
(ख) \ce{K2O} $47.10/417.25 = \qty{11.3}{\percent}$; \ce{MgO} $3 \times 40.304/417.25 =
\qty{29.0}{\percent}$; \ce{Al2O3} $50.98/417.25 = \qty{12.2}{\percent}$; \ce{SiO2} $3 \times
60.083/417.25 = \qty{43.2}{\percent}$; योग \qty{95.7}{\percent}। (ग) माइक्रोप्रोब हाइड्रोजन का
संसूचन नहीं करता: किसी अभ्रक के विश्लेषण में उसका जल, लगभग \qty{4}{\percent}, नहीं होता, और उसका योग लगभग
96 होना चाहिए।
\end{solution}

\begin{solution}{exo:b2:field-lab-methods-2:8}
(क) चरघातांकी नियम से $\beta = \ln(0.1194/0.11905)/\ln(85.909261/87.905612) = 0.002936/(-0.02297) = -0.1278$।
(ख) $0.71062 \times (86.908878/85.909261)^{-0.1278} = 0.70957$। (ग) मापा गया 86/88 सच्चे अनुपात से
नीचे है: भारी समस्थानिक को वरीयता मिली, और मापा गया 87/86 बहुत ऊँचा है।
\end{solution}

\begin{solution}{exo:b2:field-lab-methods-2:9}
(क) $0.1000 \times 15.271 \times 10^{-6}/83.9242 \times 1.3599 \times 10^{-3} = \qty{2.474e-11}{mol}$।
(ख) \eqref{eq:b2:field-lab-methods-2:id} से, $N_x = 2.474 \times 10^{-11} \times (732.8 -
0.874)/(0.874 - 0.0568) = \qty{2.216e-8}{mol}$ $^{86}\mathrm{Sr}$; स्ट्रॉन्शियम: $2.216 \times
10^{-8}/0.0986 \times 87.62 = \qty{1.969e-5}{g}$, \qty{0.0500}{g} में, \qty{394}{ppm} (393.9)।
(ग) \qty{99.95}{\percent} की पुनर्प्राप्ति: अनुपात के पूर्णांकन की सीमा में विधि शुद्ध है।
\end{solution}

\begin{solution}{exo:b2:field-lab-methods-2:10}
(क) $R_m = \sqrt{732.8 \times 0.05680} = 6.45$, जहाँ $\mathcal M = 1.018$। (ख) 0.1 पर $\mathcal M = 2.31$,
0.5 पर 1.13, 1.664 पर 1.038 और 100 पर 1.16। (ग) $R_x$ के निकट मिश्रण नमूने से मुश्किल से भिन्न होता है,
$R_s$ के निकट स्पाइक से, और सूचना वहन करने वाला अंतर $R_m$ की त्रुटि में खो जाता है:
$\mathcal M$ दोनों सिरों पर अपसरित होता है। उनके बीच न्यूनतम बहुत सपाट है
(लगभग 0.4 से 100 तक $\mathcal M < 1.2$), इसलिए विश्लेषक को स्पाइक की केवल सही कोटि-मात्रा चाहिए,
जो एक मोटी सांद्रता, उदाहरण के लिए XRF से, दे देती है।
\end{solution}

\begin{solution}{exo:b2:field-lab-methods-2:11}
(क) $\sigma = \sqrt{0.62 \times 0.38/400} = 0.0243$: $62 \pm 2.4$। (ख) $(65 - 62)/2.43 = 1.2$
मानक विचलन: चट्टान ग्रैनोडायोराइट भी हो सकती है। (ग) $\sigma = 0.015$ के लिए $0.62 \times
0.38/0.015^2 = 1047.1$ चाहिए, अतः 1048 फ़ेल्डस्पार बिंदु।
\end{solution}

\begin{solution}{exo:b2:field-lab-methods-2:12}
(क) हर बिंदु प्रायिकता $1 - p$ से चूकता है, स्वतंत्र रूप से: $(1-p)^n$। (ख) $0.99^n \le 0.05$
से $n \ge \ln 0.05/\ln 0.99 = 298.1$ मिलता है: 299 बिंदु। (ग) $n \ge \ln 0.05/\ln 0.9995 = 5990$: ग्रिड-गणना
गलत औज़ार है; किसी दुर्लभ सहायक खनिज को पूरी काट को छानकर खोजा जाता है।
\end{solution}

\begin{solution}{pb:b2:field-lab-methods-2:1}
\noindent\textbf{भाग I।}
\textbf{1.} क्वार्ट्ज़ $372/1200 = \qty{31.0}{\percent}$, क्षार-फ़ेल्डस्पार 29.0, प्लैजियोक्लेज़ 27.0, बायोटाइट
9.0, मस्कोवाइट 3.0, अपारदर्शी ऑक्साइड 1.0।
\textbf{2.} $\sqrt{0.31 \times 0.69/1200} = 0.0134$; मस्कोवाइट: $\sqrt{0.97/(0.03 \times 1200)} =
\qty{16.4}{\percent}$।
\textbf{3.} $M = 9 + 3 + 1 = \qty{13.0}{\percent}$; $M' = 13 - 3 = \qty{10.0}{\percent}$, बिना
मस्कोवाइट के: ल्यूकोक्रेटी।
\textbf{4.} $Q + A + P = 1044$ बिंदु: Q $= 35.6$, A $= 33.3$, P $= 31.0$।
\textbf{5.} $324/672 = \qty{48.2}{\percent}$, 20 और 60 के बीच Q के साथ: एक ग्रेनाइट, और 35 और
65 के बीच होने से एक मॉन्ज़ोग्रेनाइट।
\textbf{6.} $\sqrt{0.482 \times 0.518/672} = 0.0193$: 35 और 65 पर की सीमाएँ 6.9 और 8.7
मानक विचलन दूर हैं। नाम सुरक्षित है।
\textbf{7.} $0.31 \times 0.69/0.01^2 = 2139$ बिंदु।

\noindent\textbf{भाग II।}
\textbf{8.} $0.0051/0.8000 = \qty{0.64}{\percent}$।
\textbf{9.} $0.7949/8.7949 = 0.0904$, और काँच में $72 \times 0.0904 = \qty{6.51}{\percent}$ \ce{SiO2}:
संगलन नमूने को ग्यारह गुना तनु कर देता है।
\textbf{10.} $\bar I = 24.91$, $\bar c = 62.283$, $\sum(I - \bar I)^2 = 31.918$, $\sum(I - \bar I)(c -
\bar c) = 80.009$: $a = 2.5067$ प्रति kcps, $b = 62.283 - 2.5067 \times 24.91 = -0.16$।
\textbf{11.} $2.5067 \times 28.63 - 0.16 = \qty{71.61}{\percent}$।
\textbf{12.} BCR-2 $-0.04$, AGV-2 $+0.05$, RGM-1 $-0.01$: सभी विधि के मानक विचलनों (0.1 से 0.2)
से छोटे, इसलिए इस परास पर रेखा पर्याप्त है।
\textbf{13.} $71.61 + 27.15 + 0.64 = \qty{99.40}{\percent}$, 97 और 101 के बीच: दोहराया नहीं जाता।
\textbf{14.} $71.61 \times 100/(71.61 + 27.15) = \qty{72.51}{\percent}$।

\noindent\textbf{भाग III।}
\textbf{15.} $0.0986/0.8258 = 0.1194$: प्रसामान्यीकरण का मान प्राकृतिक अनुपात ही है।
\textbf{16.} घातांक $\beta = 0.007313/(-0.022972) = -0.3183$।
\textbf{17.} $0.71655 \times (86.908878/85.909261)^{-0.3183} = 0.71392$।
\textbf{18.} मापा गया 86/88 सच्चे मान से $0.11853/0.1194 - 1 = \qty{-0.73}{\percent}$ दूर है,
दो द्रव्यमान-इकाइयों पर: प्रति द्रव्यमान-इकाई लगभग \qty{0.36}{\percent}, भारी समस्थानिकों के पक्ष में।
\textbf{19.} स्पाइक का $^{86}\mathrm{Sr}/^{88}\mathrm{Sr}$ 0.7661 है: मिश्रण का 86/88 अब
प्राकृतिक 0.1194 नहीं रहा, और उसे 0.1194 पर स्थिर करना स्पाइक को उपकरणीय पूर्वाग्रह मान लेना होगा।

\noindent\textbf{भाग IV।}
\textbf{20.} 88 से अनुपातों का योग: $561.4 + 0.7661 + 0.1819 + 1 = 563.348$; $561.4/563.348 =
\qty{99.65}{\percent}$ और $0.7661/563.348 = \qty{0.13599}{\percent}$।
\textbf{21.} $0.2500 \times 15.271 \times 10^{-6}/83.9242 = \qty{4.549e-8}{mol}$ स्ट्रॉन्शियम, और
$\times 1.3599 \times 10^{-3} = \qty{6.186e-11}{mol}$ $^{86}\mathrm{Sr}$।
\textbf{22.} $561.4/0.7661 = 732.8$ और $0.0056/0.0986 = 0.05680$।
\textbf{23.} $N_x = 6.186 \times 10^{-11} \times (732.8 - 1.6812)/(1.6812 - 0.0568) =
\qty{2.784e-8}{mol}$।
\textbf{24.} $\mathcal M = 1.6812 \times 732.74/(731.12 \times 1.6244) = 1.04$; अनुपात
$1.04 \times \qty{0.01}{\percent}$ जोड़ता है, स्पाइक $0.031/15.271 = \qty{0.20}{\percent}$; वर्ग-योग में
\qty{0.20}{\percent}, जिसमें स्पाइक का अंशांकन प्रधान है।
\textbf{25.} स्ट्रॉन्शियम: $2.784 \times 10^{-8}/0.0986 = \qty{2.824e-7}{mol}$, अर्थात् $2.824 \times
10^{-7} \times 87.62 = \qty{2.474e-5}{g}$, \qty{0.1000}{g} में। वार्डन ग्रेनाइट में
\textbf{\qty[text-series-to-math=true]{247}{ppm}} स्ट्रॉन्शियम है, $\pm 0.5$।
\end{solution}
